\begin{table}[t]
\caption{Primary and exploratory QML endpoints from the frozen evidence package. All values are development-only. Lower NRMSE and Brier score are better; higher recall is better for retaining actually feasible plans. The feasibility endpoint uses the fixed 0.5 threshold.}
\label{tab:gate5-results}
\centering
\footnotesize
\begin{tabularx}{\textwidth}{@{}>{\raggedright\arraybackslash}p{0.18\textwidth}>{\raggedright\arraybackslash}p{0.29\textwidth}Y@{}}
\toprule
Stage & Model and metric & Development result and meaning \\
\midrule
Primary Gate 5 & Fidelity quantum kernel (Q01); mean NRMSE & Q01: 0.646614; physics-residual C06: 0.008739. Protocol FAIL; no predefined subgroup passed. \\
Projected follow-up & Projected quantum kernel (Q01b; P001); mean NRMSE & Q01b: 0.661207; C06: 0.006833. Exploratory negative; about 97 times the C06 error. \\
Feasibility follow-up & Feasibility-only quantum kernel (FQK; P001); recall at 0.5 & Recall: 0.108860. It retained only 10.9\% of actually feasible cases, and classical controls were stronger. \\
Future lesson & Calibrated logistic model; feasible-case recall and Brier score & Recall: 0.804253; Brier: 0.142231. Protocol-design signal, not a safety result. \\
\bottomrule
\end{tabularx}
\end{table}
